% T18 source: computational_results.tex lines 60--143
\subsection{Decomposition and Its Interaction with \texorpdfstring{$P$}{P}}

In EXP-04, the exact $T^*(m)$ profile on the finite grid was U-shaped in 12 of
the 13 curves---in every case with positive decomposition overhead. The common
zero-overhead slice
reached a plateau once serialization was removed. With low overhead, the
optimum was usually attained at $m=3$; as overhead increased, it shifted toward
$m=2$. Decomposition is therefore not monotonically beneficial: it reduces
structural serialization but adds task-specification, acceptance, and
integration overhead.

\begin{figure}[H]
  \centering
  \includegraphics[width=0.94\textwidth]{exp04_decomposition_curves.png}
  \caption{Exact decomposition curves. On the finite grid examined, positive
  overhead produces a finite optimal granularity; zero overhead produces a plateau.}
  \label{fig:exp04-decomposition}
\end{figure}

EXP-05 examines the interaction effect. Of the 100 interior interaction
contrasts with positive overhead, 96 are positive, four are zero, and none are
negative. However, the absolute gain from decomposition at $P>1$ is positive at
only 74 points and negative at 26; at 22 points, a positive interaction contrast
coincides with a negative absolute gain. Complementarity therefore means that
the effect is more favorable than at $P=1$, not that decomposition is
unconditionally beneficial. Increasing the human component of overhead
attenuates the gain. This is a result on the selected finite grid, not a
universal theorem.

\subsection{Scale and Parameter Robustness: EXP-06}

As an implementation check, EXP-06 reproduced scale invariance at all 20 of 20
exact-solver points. The advantage of illustrative Variant 4 over Variant 3 was
preserved in all 5000 random pairs and all 382 prespecified adverse pairs. This
result establishes robustness only within the frozen error families, not under
arbitrary perturbations of all parameters.

\subsection{Transfer across Synthetic DAGs: EXP-07}

Across 1440 scenarios constructed from 480 synthetic DAGs with three placements
of human work, EXP-07, in its pilot mode, supported two prespecified median
directions for \texttt{bottom\_level\_fcfs\_v1}. In C2, the median shift in
$P_{\mathrm{opt}}:=\min\operatorname*{arg\,min}_{P}T_\pi(P)$ after introducing
agent and human costs was zero; shifts in the opposite, positive direction
occurred in 0.42\% and 0.63\% of DAGs, respectively, below the prespecified
5\% threshold. In C3, the median $I_\pi(4,3)$ was 6.4, with a 95\% bootstrap
interval of $[6.0,6.8]$, while the median attenuation of the gain due to overhead
was 2.0, with an interval of $[2.0,2.0]$. Negative individual contrasts
nevertheless accounted for 6.53\% and 3.47\%, respectively. Thus, the pilot
supported aggregate directions, not universal signs for every DAG or claims
about a certified $T^*$. Because the pilot-extension rule was only partially
implemented, these decisions pertain to the resulting 20-seed pilot and require
reassessment under the fully implemented gate.

The hypotheses concerning the ``synchronous'' profile and simple concentration
of human work were not supported: the former did not, in fact, model
synchronization in time, while the latter conflated a change in total $H$ with
its placement. The retrospective interpretation of the former is therefore
limited to a comparison of equal and permuted asynchronous phase profiles, and
the conclusion regarding the placement of $H$ is retested in EXP-08. In the
EXP-07 exact audit, only 17 of 36 runs returned \texttt{OPTIMAL}; four complete
C1 pairs yielded a small estimate of the opposite sign. This audit did not
cover C2--C3 and does not replace the primary fixed-policy estimand.

\subsection{Placement of Fixed Human Work: EXP-08}

The corrected paired design of EXP-08 preserved the DAG, $A$, $H$, and total
work. With $C=\gamma=1$, moving human work between tasks changed neither
$H,W_4,L_4,B_4$ nor the active branches in any of the 2880 checks. When
$\gamma>C$, moving an amount $\Delta H$ onto the prespecified path changes the
length of that path by exactly $(\gamma-C)\Delta H$; this identity was confirmed
in all 480 pairs. The change in the fixed-policy makespan, however,
depended on the phases and on $P$. For the primary \texttt{io} profile, the
normalized paired median ``high minus low'' at $P=4$ was $-0.0077$, with a 95\%
bootstrap interval of $[-0.0090,-0.0065]$, and was classified as practically
equivalent under the prespecified margin of 0.01. At $P=8$, the estimate was
$-0.0117$, with an interval of $[-0.0120,-0.0101]$, and the decision was
negative.

All 16 exact-audit scenarios had status \texttt{OPTIMAL}, but this does not make
the statistical conclusion determinate: the eight certified pairs yielded the
same median, $0.0010$, for $P=4$ and $P=8$, with an interval of
$[-0.0150,0.0080]$, so both decisions remained inconclusive. Thus, two distinct
characteristics matter: total human work determines the resource ceiling; when
$\gamma\ne C$, its placement changes the path weights and $L_4$, whereas when
$C=\gamma$, phase order can still change the queue and the fixed-policy makespan
while the aggregates remain unchanged.
